\section{Converting between line descriptions}

The same line can be written in many forms. Therefore one of the basic skills in analytic geometry is translation. We translate from the form that is given to the form that makes the next question easier.

This translation should not be mechanical. Each conversion has a meaning. When we convert a parametric equation into an implicit equation, we are changing from a direction-based description to a condition-based description. When we convert an implicit equation into a parametric equation, we are finding a point and a direction hidden inside a condition.

\subsection*{From parametric form to implicit form}

Suppose a line is given in parametric form:
\[
P=P_0+t v.
\]
In the plane, if
\[
P_0=(x_0,y_0),
\qquad
v=(a,b),
\]
then a normal vector may be chosen as
\[
n=(-b,a).
\]
The implicit equation is then obtained from
\[
n\cdot(P-P_0)=0.
\]

\begin{examplebox}
Convert the line
\[
(x,y)=(1,2)+t(3,-1)
\]
to implicit form.

The direction vector is
\[
v=(3,-1).
\]
A normal vector is
\[
n=(1,3),
\]
because
\[
(3,-1)\cdot(1,3)=3-3=0.
\]
Using the point \((1,2)\), the normal form is
\[
1(x-1)+3(y-2)=0.
\]
Thus
\[
x-1+3y-6=0,
\]
so
\[
x+3y-7=0.
\]
\end{examplebox}

The direction vector was visible in the parametric form. The normal vector became visible after conversion to implicit form.

\subsection*{From implicit form to parametric form}

Now suppose the line is given by
\[
Ax+By+C=0.
\]
A normal vector is immediately visible:
\[
n=(A,B).
\]
A direction vector must be perpendicular to \(n\). One convenient choice is
\[
v=(-B,A),
\]
because
\[
(A,B)\cdot(-B,A)=-AB+AB=0.
\]
We also need one point on the line. This can be found by choosing one coordinate and solving for the other, whenever possible.

\begin{examplebox}
Convert
\[
2x-5y+10=0
\]
to parametric form.

A normal vector is
\[
n=(2,-5).
\]
A direction vector is
\[
v=(5,2),
\]
because
\[
(2,-5)\cdot(5,2)=10-10=0.
\]
To find a point, choose \(x=0\). Then
\[
-5y+10=0,
\]
so
\[
y=2.
\]
Thus \((0,2)\) lies on the line. A parametric equation is
\[
(x,y)=(0,2)+t(5,2).
\]
\end{examplebox}

There are many possible choices of the point. If we choose another point on the same line, we get another parametrization of the same geometric object.

\subsection*{From two points to implicit form}

If a line is given by two distinct points, first build a direction vector. Then choose a normal vector perpendicular to that direction vector.

\begin{examplebox}
Find an implicit equation of the line through
\[
A=(1,-1),
\qquad
B=(4,5).
\]
The direction vector is
\[
B-A=(3,6).
\]
We may simplify the direction to
\[
v=(1,2).
\]
A normal vector is
\[
n=(-2,1),
\]
because
\[
(1,2)\cdot(-2,1)=-2+2=0.
\]
Using the point \(A=(1,-1)\), the normal form is
\[
-2(x-1)+1(y-(-1))=0.
\]
Thus
\[
-2(x-1)+y+1=0.
\]
Expanding,
\[
-2x+2+y+1=0,
\]
so
\[
-2x+y+3=0.
\]
Equivalently,
\[
2x-y-3=0.
\]
\end{examplebox}

Both final equations describe the same line. Multiplying every coefficient by a non-zero number does not change the set of solutions.

\subsection*{From slope form to direction and normal vectors}

If a non-vertical line has equation
\[
y=mx+b,
\]
then a direction vector is
\[
v=(1,m).
\]
Indeed, when \(x\) increases by \(1\), the value of \(y\) changes by \(m\).

After rewriting as
\[
mx-y+b=0,
\]
we also see a normal vector:
\[
n=(m,-1).
\]
These two vectors are perpendicular, because
\[
(1,m)\cdot(m,-1)=m-m=0.
\]

\begin{examplebox}
For the line
\[
y=-3x+4,
\]
a direction vector is
\[
v=(1,-3).
\]
Rewriting as
\[
3x+y-4=0,
\]
we see a normal vector
\[
n=(3,1).
\]
Indeed,
\[
(1,-3)\cdot(3,1)=3-3=0.
\]
\end{examplebox}

\subsection*{Checking the translation}

After converting one form into another, it is good practice to check that the descriptions agree. We can check that the chosen point satisfies the new equation. We can check that the direction vector is perpendicular to the normal vector. We can also test a second point generated by the parametrization.

\begin{summarybox}
When converting between descriptions, ask:
\begin{itemize}
\item What information is visible in the given form?
\item Do I need a direction vector or a normal vector?
\item Have I found at least one point on the line?
\item Is the direction vector perpendicular to the normal vector?
\item Does the converted equation describe the same set of points?
\end{itemize}
\end{summarybox}

Conversion is not a ritual of rearranging symbols. It is a way of making hidden geometric information visible.